\documentclass{beamer}
\usetheme{Madrid}

\title{Introduction to Beamer}
\author{Umed harwani}
\institute{MNIT Jaipur}
\date{\today}

\begin{document}

\begin{frame}
    \titlepage
\end{frame}

\begin{frame}{What is Beamer?}
Beamer is a LaTeX class for making slides.
You write the slides as code and compile them to a PDF.
\end{frame}

\begin{frame}{Why I Use It}
\begin{itemize}
    \item<1-> Equations look clean
    \item<2-> Layout stays consistent across slides
    \item<3-> Overlays reveal points one at a time
\end{itemize}
\end{frame}

\begin{frame}{Steps to Make a Deck}
\begin{enumerate}
    \item Pick the \texttt{beamer} class
    \item Fill in title, author, date
    \item Add one \texttt{frame} per slide
    \item Compile with \texttt{pdflatex}
\end{enumerate}
\end{frame}

\begin{frame}{Equations}
Quadratic formula:
\[
x=\frac{-b\pm\sqrt{b^2-4ac}}{2a}
\]
\end{frame}

\begin{frame}[fragile]{Code on a Slide}
\begin{verbatim}
#include <iostream>
int main() {
    std::cout << "Hello World";
    return 0;
}
\end{verbatim}
\end{frame}

\begin{frame}{Summary}
Beamer is a good fit for technical talks: equations, code and
step-by-step reveals all work out of the box.
\end{frame}

\end{document}
